\documentclass[a4paper,11pt]{article}

\usepackage[czech,english]{babel}
% Fonts %
\usepackage{fouriernc}
\usepackage[T1]{fontenc}

\usepackage{advdate}

% Colors %
\usepackage[dvipsnames]{color}
\usepackage{xcolor}

% Page Layout %
\usepackage[margin=1in]{geometry}
\usepackage{paracol}

% Fancy Headers %
\usepackage{fancyhdr}
\fancyhf{}
\rhead{\sffamily\large Vocabulary Exam -- 2.AB PreIB Math}
\lhead{\sffamily\large \today}
\renewcommand{\headrulewidth}{0.4pt}
\setlength{\headheight}{16pt}

% Math
\usepackage{mathtools}
\usepackage{amssymb}
\usepackage{faktor}
\usepackage{import}
\usepackage{caption}
\usepackage{subcaption}
\usepackage{wrapfig}
\usepackage{enumitem}
\usepackage{tikz}
\usetikzlibrary{cd,positioning,babel,shapes}
\usepackage{tkz-base}
\usepackage{tkz-euclide}

% Theorems
\usepackage{thmtools}
\usepackage[thmmarks, amsmath, thref]{ntheorem}

% Title %
\title{\Huge\textsf{Vocabulary Exam -- 2.AB PreIB Math}\\
  \Large\textsf{Logic \& Set Theory}
  \author{Áďa Klepáčů}
  \date{\today}
}

% Table of Contents %
\usepackage{hyperref}
\hypersetup{
  colorlinks=true,
  linktoc=all,
  linkcolor=blue
}

% Tables %
\usepackage{booktabs}
\usepackage{tabularx}

% Patch for hyphens
\usepackage{regexpatch}
\makeatletter
% Change the `-` delimiter to an active character
\xpatchparametertext\@@@cmidrule{-}{\cA-}{}{}
\xpatchparametertext\@cline{-}{\cA-}{}{}
\makeatother

\newcolumntype{s}{>{\centering\arraybackslash}p{.4\textwidth}}

% Operators %
\DeclareMathOperator{\Ker}{Ker}
\DeclareMathOperator{\Img}{Im}
\DeclareMathOperator{\End}{End}
\DeclareMathOperator{\Aut}{Aut}
\DeclareMathOperator{\Inn}{Inn}

% Common operators %
\newcommand{\R}{\mathbb{R}}
\newcommand{\N}{\mathbb{N}}
\newcommand{\Z}{\mathbb{Z}}
\newcommand{\Q}{\mathbb{Q}}
\newcommand{\C}{\mathbb{C}}

\newcommand{\tr}{\textcolor{red}}
\newcommand{\tb}{\textcolor{blue}}
\newcommand{\tg}{\textcolor{green}}
\newcommand{\tm}{\textcolor{magenta}}
\newcommand{\tv}{\textcolor{violet}}

% American Paragraph Skip %
\setlength{\parindent}{0pt}
\setlength{\parskip}{1em}

% Document %
\pagestyle{fancy}
\renewcommand{\baselinestretch}{1.2}
\newcommand{\blank}{\underline{\hspace{12ex}}}
\begin{document}

\thispagestyle{fancy}


\columnratio{.7}
\begin{paracol}{2}
  In traditional mathematics, we distinguish a few different number sets. The
  basis for all of them is the set of numbers we use for `counting' things --
  these numbers are called \blank. Operations on these numbers include \blank,
  \blank~and \blank~with the next operation being basically just a convenient
  way to express repetition of the previous one.
  
  Next in line is the set of \blank~-- numbers which can be positive or negative
  and thus moreover express notions like `absence' or `debt'. Often defined as
  \blank~between two positive numbers, they also include the operation of
  \blank~which isn't sensible on positive numbers as the \blank~of two positive
  numbers may be negative.

  If we also in need of numbers able to express \emph{parts of objects},
  \blank~numbers come to the rescue. We write them as either decimal numbers or
  \blank, with the latter form getting across the idea that they are essentially
  results of dividing one whole number by another: or, as we call the results of
  division, \blank.

  Finally, there are also irrational numbers -- numbers which literally can't be
  written as a `ratio' of two whole numbers. They come about for example as
  \blank~of many positive numbers that are not powers of other numbers, like
  $2$, $3$ or $5$. The numbers that also include irrational numbers and are
  heavily used in geometry are called \blank.

  \switchcolumn
  \renewcommand{\arraystretch}{2}
  \vspace*{-32pt}
  \begin{tabular}{c}
    \textbf{EXPONENTIATION}\\
    \textbf{SUBTRACTION}\\
    \textbf{INTEGER}\\
    \textbf{QUOTIENT}\\
    \textbf{REAL}\\
    \textbf{FRACTION}\\
    \textbf{MULTIPLICATION}\\
    \textbf{ROOT}\\
    \textbf{RATIONAL}\\
    \textbf{ADDITION}\\
    \textbf{DIFFERENCE}\\
    \textbf{NATURAL}\\
  \end{tabular}
\end{paracol}

\end{document}
